\chapter{Introducción}
\label{cap:introduccion}


\section{Antecedentes}
\label{sec:introduccion:antecedentes}

La medicina nuclear ha experimentado en las últimas décadas una transición hacia la medicina de precisión, impulsada en gran medida
por el auge de la \textbf{teranosis} (acrónimo de terapia y diagnóstico). Este enfoque se basa en el uso de pares de isótopos radiactivos
("matched pairs") que comparten propiedades químicas idénticas o muy similares, permitiendo utilizar uno para el diagnóstico por
imagen y otro para la terapia en la misma diana molecular.

Uno de los pares teranósticos más relevantes en la práctica clínica actual es el formado por el \textbf{Itrio-86} (\yochoseis) y el
\textbf{Itrio-90} (\ynuevecero). El \ynuevecero es un emisor beta negativo puro ($\beta^-$). Sin embargo, su emisión carece de
fotones gamma detectables, dificulta su detección mediante las técnicas de imagen convencionales, imposibilitando una dosimetría
precisa previa al tratamiento. Su isótopo análogo, el \yochoseis es un emisor de positrones ($\beta^+$) que permite visualizar la
biodistribución del fármaco mediante Tomografía por Emisión de Positrones (PET).

A pesar de su utilidad, la imagen cuantitativa con \yochoseis presenta desafíos físicos significativos. Se clasifica como un
\textbf{isótopo "sucio"}, ya que ya que su desintegración conlleva la emisión simultánea de una cascada de rayos gamma de alta
energía (\textit{prompt gammas}). La literatura científica (\cite{Lubberink2011,Conti2016}) señala que estos gammas generan
coincidencias triples y aumentan la tasa de eventos aleatorios, lo que los escáneres comerciales estándar malinterpretan,
degradando severamente el contraste y la exactitud cuantitativa (\cite{Buchholz2003}).

Para abordar esta problemática desde una perspectiva industrial y tecnológica, la simulación Monte Carlo es la herramienta
fundamental. En este trabajo se utiliza el framework \textbf{\penred} (\cite{Gimnez-Alventosa2021})
basado en PENELOPE, para modelar con alta fidelidad la geometría y la física de detección del escáner preclínico
\textbf{Bruker PET/CT Si78}. El uso de las herramientas de reconstrucción propietarias de la empresa (Bruker RecoServer)
permite validar los resultados en un entorno de aplicación real.

\section{Motivación}
\label{sec:introduccion:motivacion}

La motivación principal de este trabajo surge de una necesidad clínica y tecnológica no resuelta: \textbf{la falta de precisión
    cuantitativa} en las imágenes PET de \yochoseis impide aprovechar todo el potencial del par teranóstico \yochoseis/\ynuevecero.

Este proyecto se desarrolla en el marco de una colaboración con la empresa \textbf{Bruker}, líder mundial en instrumentación científica y
diagnóstico por imagen. La motivación técnica reside en optimizar el rendimiento de sus sistemas preclínicos para isótopos complejos.
Mediante el uso de simulación computacional avanzada, se pretende descomponer los efectos degradantes (\emph{scatter}, atenuación y \emph{prompt gammas})
y proponer una \textbf{optimización} de los parámetros de \textbf{adquisición y procesado} (ventanas de energía y tiempo) específica para sus detectores.
Se busca demostrar que, mediante un ajuste fino, es posible recuperar el contraste perdido y transformar una imagen ruidosa en una
herramienta dosimétrica fiable.

Además, en este trabajo se explora una motivación innovadora que va más allá de la corrección: el \textbf{cambio de paradigma} hacia la imagen
\textbf{multiplexada}. Se plantea la base para utilizar los \emph{prompt gammas} no como ruido, sino como una
"firma"\space para distinguir el \yochoseis de otros isótopos, abriendo la puerta a adquisiciones simultáneas duales (mPET) (\cite{Pratt2023}).

\section{Carácter Académico del Trabajo}
\label{sec:introduccion:caracter}

Desde el punto de vista académico, la realización de este Trabajo de Fin de Grado (TFG) tiene como objeto la integración y
puesta en práctica de las competencias adquiridas a lo largo del Grado en Ingeniería Biomédica. Este proyecto permite demostrar
la capacidad para:

\begin{itemize}
    \item Aplicar fundamentos de física de radiaciones y tecnología de imagen médica a problemas reales de la industria tecnológica sanitaria.
    \item Utilizar herramientas de ingeniería computacional (simulación Monte Carlo, C++, Python) en un entorno de desarrollo profesional.
    \item Diseñar una metodología experimental rigurosa para la validación de equipos médicos de alta tecnología (\textbf{Bruker Si78}).
    \item Elaborar documentación técnica y científica siguiendo los estándares de la ingeniería.
\end{itemize}

En definitiva, este trabajo constituye un ejercicio de ingeniería aplicada a la salud, donde se combinan conocimientos
multidisciplinares para proponer una mejora tecnológica con impacto directo en la práctica clínica y la investigación biomédica.

\section{Alcance del proyecto}
\label{sec:introduccion:alcance}

El presente Trabajo de Fin de Grado abarca el ciclo completo de investigación in silico, desde el modelado físico hasta
el análisis de imagen médica, centrándose en la problemática específica del isótopo Itrio-86 (\yochoseis) en entornos
preclínicos.

El alcance del trabajo se estructura en los siguientes bloques principales:

\subsection{Modelado y Simulación Computacional}
El proyecto comprende la configuración y ejecución de simulaciones Monte Carlo utilizando el framework \textbf{penRed}
(\cite{Gimnez-Alventosa2021}).

\begin{itemize}
    \item Se incluye el modelado geométrico preciso mediante el software
          \textbf{Blender} del escáner preclínico comercial Bruker PET/CT Si78 (\cite{unknown-author-no-date}).
    \item Se abarca la simulación de la física de desintegración compleja del \yochoseis (emisión de positrones y cascadas de prompt gammas)
          mediante el módulo PENNUC, así como la simulación de un isótopo de referencia (\fdiezocho) para estudios comparativos.
    \item El escenario de simulación se limita al uso del fantoma éstandar de calidad de imagen (Derenzo), permitiendo estudios
          reproducibles de resolución y contraste.
\end{itemize}

\subsection{Desarrollo de software de Post-procesado}

El alcance comprende la implementación y validación de algoritmos de procesamiento de datos
crudos (singles) para extraer las coincidencias PET correspondientes, en formato \emph{list mode} (LM).
\footnote{Cabe aclarar que los códigos de Post-procesado iniciales, los cuales no contenían las funcinones
    en relación al multiplexado me fueron proporcionados por Vicent. Partiendo de esa base,
    yo los modifiqué para que permiteran multiplexing. Yo los he completado, mas no los desarrollé desde el principio}


Se incluye el diseño de un algoritmo de clasificación de eventos para la clasificación
de eventos triples basado en la tecnica de separación
de isótopos de \cite{Andreyev2011}, implementado en C++.

\subsection{Estudio de Optimización de Parámetros}

El trabajo incluye un análisis cuantitativo para determinar los parámetros de adquisición óptimos.
El alcance de este estudio incluye:
\begin{itemize}
    \item Un barrido sistemático de Ventanas de Energía (EW) y Ventanas de Coincidencia (CW).
    \item La evaluación de la calidad de imagen mediante métricas estandarizadas: Coeficiente de Recuperación de Contraste (CRC) y Relación Señal-Ruido (SNR).
    \item La reconstrucción de imágenes utilizando el software comercial Bruker RecoServer bajo el algoritmo estándar MLEM.
\end{itemize}
\footnote{esto lo hice previo a desarrollar el algoritmo de separación de isótopos, realmente}

\subsection{Prueba de Concepto de Imagen Multiplexada}

Finalmente, el proyecto alcanza la implementación de una prueba de concepto para la técnica de PET
mutliplexado (mPET).

Se modeló en Blender, un fantoma basado en el de \cite{Andreyev2011}, el cual contenía diferentes
varillas con diferente concentración de dos isótopos, uno sucio (\yochoseis) y otro limpio (\fdiezocho).

Se realizaron simulaciones con penRed, simulando altas energías para incluir todos los 
gammas. Luego se usó el algoritmo de postProcess modificado con el nuevo metodo mutliplexing,
aprovechando estos gammas de alta energía , para generar los distintos datasets, y finalmente se 
usaron unos scripts de python para leer las imágenes crudas tras el proceso de
reconstrucción y sustraer la imagen ($A+B - F*B'$). 

\section{Estructura de la memoria}

% Terminar más adelante cuando tenga la memoria hecha??


\section{Repositorio}

El código fuente desarrollado para el post-procesado y los scripts de análisis están 
disponibles en: https://github.com/gutierrels/dirtyIsotopes/. Más iformación en \ref{fig:qr_repo}



¿Qué SÍ subir?
\begin{itemize}
    \item Tus scripts de C++ (createCoincidences.cpp, etc.).
    \item Los archivos de configuración de penRed (config.in) genéricos.
    \item Los scripts de Python para generar gráficas.
    \item El archivo .geo del fantoma Derenzo.
\end{itemize}

